% GENERATED FILE. Edit canonical lessons, then run npm run generate:books.
% canonical-source-hash: fnv1a64:62b3c56ab57aadbf
% canonical-lessons: BN-C168-shashuri, BN-C168-bor, BN-C168-kone, BN-C168-biye, BN-C168-atmiyo

\chapter{Mother-in-law, Bridegroom, Bride, Wedding, Relative}
\label{ch:bn-mother-in-law-bridegroom-bride-wedding-relative}
\hlchaptermodality{168}

\begin{chapteropening}
\textbf{By the end of this chapter:} \emph{I can say a mother-in-law, a bridegroom, a bride, a wedding and a relative.}

Complete the last lesson of chapter 168: I can say a mother-in-law, a bridegroom, a bride, a wedding and a relative.
\end{chapteropening}

\section[\texorpdfstring{\bn{শাশুড়ি} (shāshuṛi)}{শাশুড়ি (shāshuṛi)}]{\bn{শাশুড়ি} (shāshuṛi) — a mother-in-law}

\label{lesson:BN-C168-shashuri}



\begin{quote}
Before the new one: say the Bengali for a motorbike, then the Bengali for a lorry.
\end{quote}

\subsection*{You'll want to know: \bn{শাশুড়ি}}
\textbf{\bn{শাশুড়ি}} — \emph{shāshuṛi} — \textquotedblleft{}a mother-in-law\textquotedblright{}.

\begin{cousinweb}
\textbf{\bn{শাশুড়ি}} (\emph{shāshuṛi}) is a tadbhava, a word inherited from Sanskrit and reshaped by sound change. It comes from \emph{śvaśrū}, \textquotedblleft{}husband's or wife's mother\textquotedblright{} (CDIAL 12759). CDIAL lists Bengali \emph{sāsuṛi} and the shorter \emph{sās}.
\end{cousinweb}

\begin{grammarlens}[title={What you've built}]
One more word for this chapter.
\end{grammarlens}

\subsection*{Your turn}
\begin{itemize}
\raggedright
  \item \emph{Say it:} \emph{shāshuṛi}
  \item \emph{Say it:} \emph{shāshuṛi}, once more
  \item \emph{Say it:} say \emph{shāshuṛi}
  \item \emph{Recall:} say the Bengali for a motorbike, then the Bengali for a lorry, then say \emph{shāshuṛi} again
\end{itemize}

\begin{tcolorbox}[breakable,colback=teal!4,colframe=teal!35!black,title={Before you move on}]
What does \bn{শাশুড়ি} mean? (\textquotedblleft{}A mother-in-law\textquotedblright{}.) Say it once more.
\end{tcolorbox}
\section[\texorpdfstring{\bn{বর} (bôr)}{বর (bôr)}]{\bn{বর} (bôr) — a bridegroom}

\label{lesson:BN-C168-bor}



\begin{quote}
Before the new one: say the Bengali for a lorry, then the Bengali for a mother-in-law.
\end{quote}

\subsection*{You'll want to know: \bn{বর}}
\textbf{\bn{বর}} — \emph{bôr} — \textquotedblleft{}a bridegroom\textquotedblright{}.

\begin{cousinweb}
\textbf{\bn{বর}} (\emph{bôr}) is inherited from Sanskrit \emph{vara}, \textquotedblleft{}suitor, bridegroom\textquotedblright{}, and has come down with almost no change (CDIAL 11309).
\end{cousinweb}

\begin{grammarlens}[title={What you've built}]
One more word for this chapter.
\end{grammarlens}

\subsection*{Your turn}
\begin{itemize}
\raggedright
  \item \emph{Say it:} \emph{bôr}
  \item \emph{Say it:} \emph{bôr}, once more
  \item \emph{Say it:} say \emph{bôr}
  \item \emph{Recall:} say the Bengali for a lorry, then the Bengali for a mother-in-law, then say \emph{bôr} again
\end{itemize}

\begin{tcolorbox}[breakable,colback=teal!4,colframe=teal!35!black,title={Before you move on}]
What does \bn{বর} mean? (\textquotedblleft{}A bridegroom\textquotedblright{}.) Say it once more.
\end{tcolorbox}
\section[\texorpdfstring{\bn{কনে} (kône)}{কনে (kône)}]{\bn{কনে} (kône) — a bride}

\label{lesson:BN-C168-kone}



\begin{quote}
Before the new one: say the Bengali for a mother-in-law, then the Bengali for a bridegroom.
\end{quote}

\subsection*{You'll want to know: \bn{কনে}}
\textbf{\bn{কনে}} — \emph{kône} — \textquotedblleft{}a bride\textquotedblright{}.

\begin{cousinweb}
\textbf{\bn{কনে}} (\emph{kône}) appears to go back to Sanskrit \emph{kanyā}, \textquotedblleft{}maiden\textquotedblright{}, but CDIAL gives no Bengali form under that entry, only Assamese \emph{kanāi} in a ritual sense (CDIAL 2737).
\end{cousinweb}

\begin{grammarlens}[title={What you've built}]
One more word for this chapter.
\end{grammarlens}

\subsection*{Your turn}
\begin{itemize}
\raggedright
  \item \emph{Say it:} \emph{kône}
  \item \emph{Say it:} \emph{kône}, once more
  \item \emph{Say it:} say \emph{kône}
  \item \emph{Recall:} say the Bengali for a mother-in-law, then the Bengali for a bridegroom, then say \emph{kône} again
\end{itemize}

\begin{tcolorbox}[breakable,colback=teal!4,colframe=teal!35!black,title={Before you move on}]
What does \bn{কনে} mean? (\textquotedblleft{}A bride\textquotedblright{}.) Say it once more.
\end{tcolorbox}
\section[\texorpdfstring{\bn{বিয়ে} (biye)}{বিয়ে (biye)}]{\bn{বিয়ে} (biye) — a wedding}

\label{lesson:BN-C168-biye}



\begin{quote}
Before the new one: say the Bengali for a bridegroom, then the Bengali for a bride.
\end{quote}

\subsection*{You'll want to know: \bn{বিয়ে}}
\textbf{\bn{বিয়ে}} — \emph{biye} — \textquotedblleft{}a wedding\textquotedblright{}.

\begin{cousinweb}
\textbf{\bn{বিয়ে}} (\emph{biye}) is a tadbhava, a word inherited from Sanskrit and reshaped by sound change. It comes from \emph{vivāha}, \textquotedblleft{}marriage\textquotedblright{} (CDIAL 11920). CDIAL lists Bengali \emph{biyā}, \textquotedblleft{}marriage\textquotedblright{}.
\end{cousinweb}

\begin{grammarlens}[title={What you've built}]
One more word for this chapter.
\end{grammarlens}

\subsection*{Your turn}
\begin{itemize}
\raggedright
  \item \emph{Say it:} \emph{biye}
  \item \emph{Say it:} \emph{biye}, once more
  \item \emph{Say it:} say \emph{biye}
  \item \emph{Recall:} say the Bengali for a bridegroom, then the Bengali for a bride, then say \emph{biye} again
\end{itemize}

\begin{tcolorbox}[breakable,colback=teal!4,colframe=teal!35!black,title={Before you move on}]
What does \bn{বিয়ে} mean? (\textquotedblleft{}A wedding\textquotedblright{}.) Say it once more.
\end{tcolorbox}
\section[\texorpdfstring{\bn{আত্মীয়} (ātmīyô)}{আত্মীয় (ātmīyô)}]{\bn{আত্মীয়} (ātmīyô) — a relative}

\label{lesson:BN-C168-atmiyo}



\begin{quote}
Before the new one: say the Bengali for a bride, then the Bengali for a wedding.
\end{quote}

\subsection*{You'll want to know: \bn{আত্মীয়}}
\textbf{\bn{আত্মীয়}} — \emph{ātmīyô} — \textquotedblleft{}a relative\textquotedblright{}.

\begin{cousinweb}
\textbf{\bn{আত্মীয়}} (\emph{ātmīyô}) is Sanskrit \emph{ātmīya}, \textquotedblleft{}one's own\textquotedblright{}, taken over whole as a learned loan, a tatsama.
\end{cousinweb}

\begin{grammarlens}[title={What you've built}]
One more word for this chapter.
\end{grammarlens}

\subsection*{Your turn}
\begin{itemize}
\raggedright
  \item \emph{Say it:} \emph{ātmīyô}
  \item \emph{Say it:} \emph{ātmīyô}, once more
  \item \emph{Say it:} say \emph{ātmīyô}
  \item \emph{Recall:} say the Bengali for a bride, then the Bengali for a wedding, then say \emph{ātmīyô} again
\end{itemize}

\begin{tcolorbox}[breakable,colback=teal!4,colframe=teal!35!black,title={Before you move on}]
What does \bn{আত্মীয়} mean? (\textquotedblleft{}A relative\textquotedblright{}.) Say it once more.
\end{tcolorbox}
